\begin{textsample}{POS Dim 1 – 2011 – Score 20.00 – 2011/t002528.txt}  \label{ex:f1_pos_034}
I ' m a neuroscientist.
And in neuroscience, we have to deal with many difficult questions about \textbf{the} brain.
But I want to start with \textbf{the} easiest question and \textbf{the} question you really should have all asked yourselves at some point in your life, because it ' s a fundamental question if we want to understand brain function.
And that is, why do we and other animals have brains?
Not all species on our \textbf{planet} have brains, so if we want to know what \textbf{the} brain is for, let ' s think about why we \textbf{evolved} one.
Now you may reason that we have one to perceive \textbf{the} world or to think, and that ' s completely wrong.
If you think about this question for any length of time, it ' s blindingly obvious why we have a brain.
We have a brain for one reason and one reason only, and that ' s to produce adaptable and complex movements.
There is no other reason to have a brain.
Think about it.
Movement is \textbf{the} only way you have of affecting \textbf{the} world around you.
Now that ' s not quite true.
There ' s one other way, and that ' s through sweating.
But apart from that, everything else goes through contractions of muscles.
So think about communication — speech, gestures, writing, sign language — they ' re all mediated through contractions of your muscles.
So it ' s really important to remember that sensory, memory and cognitive processes are all important, but they ' re only important to either drive or suppress future movements.
There can be no \textbf{evolutionary} advantage to laying down memories of childhood or perceiving \textbf{the} color of a rose if it doesn ' t affect \textbf{the} way you ' re going to move later in life.
Now for those who don ' t believe this argument, we have trees and grass on our \textbf{planet} without \textbf{the} brain, but \textbf{the} clinching evidence is this animal here — \textbf{the} humble \textbf{sea} squirt.
Rudimentary animal, has a nervous system, swims around in \textbf{the} ocean in its juvenile life.
And at some point of its life, it implants on a \textbf{rock}.
And \textbf{the} first thing it does in implanting on that \textbf{rock}, which it never leaves, is to digest its own brain and nervous system for food.
So once you don ' t need to move, you don ' t need \textbf{the} luxury of that brain.
And this animal is often taken as an analogy to what happens at universities when professors get tenure, but that ' s a different subject.
So I am a movement chauvinist.
I believe movement is \textbf{the} most important function of \textbf{the} brain — don ' t let anyone tell you that it ' s not true.
Now if movement is so important, how well are we doing understanding how \textbf{the} brain controls movement?
And \textbf{the} answer is we ' re doing extremely poorly ; it ' s a very hard problem.
But we can look at how well we ' re doing by thinking about how well we ' re doing building machines which can do what humans can do.
Think about \textbf{the} game of chess.
How well are we doing determining what piece to move where?
If you pit Garry Kasparov here, when he ' s not in jail, against IBM ' s Deep Blue, well \textbf{the} answer is IBM ' s Deep Blue will occasionally win.
And I think if IBM ' s Deep Blue played anyone in this room, it would win every time.
That problem is solved.
What about \textbf{the} problem of picking up a chess piece, dexterously \textbf{manipulating} it and putting it back down on \textbf{the} board?
If you put a five \textbf{year-old} child ' s dexterity against \textbf{the} best robots of today, \textbf{the} answer is simple : \textbf{the} child wins easily.
There ' s no competition at all.
Now why is that top problem so easy and \textbf{the} bottom problem so hard?
One reason is a very smart five \textbf{year-old} could tell you \textbf{the} algorithm for that top problem — look at all possible moves to \textbf{the} end of \textbf{the} game and choose \textbf{the} one that makes you win.
So it ' s a very simple algorithm.
Now of course there are other moves, but with vast \textbf{computers} we approximate and come close to \textbf{the} optimal solution.
When it comes to being dexterous, it ' s not even clear what \textbf{the} algorithm is you have to solve to be dexterous.
And we ' ll see you have to both perceive and act on \textbf{the} world, which has a lot of problems.
But let me show you cutting-edge robotics.
Now a lot of robotics is very impressive, but manipulation robotics is really just in \textbf{the} dark ages.
So this is \textbf{the} end of a Ph.
D. project from one of \textbf{the} best robotics institutes.
And \textbf{the} student has trained this robot to pour this water into a glass.
It ' s a hard problem because \textbf{the} water sloshes about, but it can do it.
But it doesn ' t do it with anything like \textbf{the} agility of a human.
Now if you want this robot to do a different task, that ' s another three-year Ph.
D. program.
There is no generalization at all from one task to another in robotics.
Now we can compare this to cutting-edge human performance.
So what I ' m going to show you is Emily Fox winning \textbf{the} world record for cup \textbf{stacking}.
Now \textbf{the} Americans in \textbf{the} audience will know all about cup \textbf{stacking}.
It ' s a high school sport where you have 12 cups you have to \textbf{stack} and unstack against \textbf{the} clock in a prescribed order.
And this is her getting \textbf{the} world record in real time.
And she ' s pretty happy.
We have no idea what is going on inside her brain when she does that, and that ' s what we ' d like to know.
So in my group, what we try to do is reverse engineer how humans control movement.
And it sounds like an easy problem.
You send a command down, it causes muscles to contract.
Your arm or body moves, and you get sensory feedback from vision, from skin, from muscles and so on. \textbf{The} trouble is these signals are not \textbf{the} beautiful signals you want them to be.
So one thing that makes controlling movement difficult is, for example, sensory feedback is extremely noisy.
Now by noise, I do not mean sound.
We use it in \textbf{the} engineering and neuroscience sense meaning a random noise corrupting a signal.
So \textbf{the} old days before digital radio when you were tuning in your radio and you heard"crrcckkk"on \textbf{the} station you wanted to hear, that was \textbf{the} noise.
But more generally, this noise is something that corrupts \textbf{the} signal.
So for example, if you put your hand under a table and try to localize it with your other hand, you can be off by several centimeters due to \textbf{the} noise in sensory feedback.
Similarly, when you put motor output on movement output, it ' s extremely noisy.
Forget about trying to hit \textbf{the} bull ' s eye in darts, just aim for \textbf{the} same spot over and over again.
You have a huge spread due to movement variability.
And more than that, \textbf{the} outside world, or task, is both ambiguous and variable. \textbf{The} teapot could be full, it could be empty.
It changes over time.
So we work in a whole sensory movement task \textbf{soup} of noise.
Now this noise is so great that society places a huge premium on those of us who can reduce \textbf{the} consequences of noise.
So if you ' re lucky enough to be able to knock a small white ball into a hole several hundred yards away using a long metal stick, our society will be willing to reward you with hundreds of millions of dollars.
Now what I want to convince you of is \textbf{the} brain also goes through a lot of effort to reduce \textbf{the} negative consequences of this sort of noise and variability.
And to do that, I ' m going to tell you about a framework which is very popular in statistics and machine \textbf{learning} of \textbf{the} last 50 years called Bayesian decision theory.
And it ' s more recently a unifying way to think about how \textbf{the} brain deals with uncertainty.
And \textbf{the} fundamental idea is you want to make inferences and then take actions.
So let ' s think about \textbf{the} inference.
You want to generate beliefs about \textbf{the} world.
So what are beliefs?
Beliefs could be : where are my arms in space?
Am I looking at a cat or a fox?
But we ' re going to represent beliefs with \textbf{probabilities}.
So we ' re going to represent a belief with a number between zero and one — zero meaning I don ' t believe it at all, one means I ' m absolutely certain.
And numbers in between give you \textbf{the} gray levels of uncertainty.
And \textbf{the} key idea to Bayesian inference is you have two sources of information from which to make your inference.
You have data, and data in neuroscience is sensory input.
So I have sensory input, which I can take in to make beliefs.
But there ' s another source of information, and that ' s effectively prior knowledge.
You accumulate knowledge throughout your life in memories.
And \textbf{the} point about Bayesian decision theory is it gives you \textbf{the} mathematics of \textbf{the} optimal way to combine your prior knowledge with your sensory evidence to generate new beliefs.
And I ' ve put \textbf{the} formula up there.
I ' m not going to explain what that formula is, but it ' s very beautiful.
And it has real beauty and real explanatory power.
And what it really says, and what you want to estimate, is \textbf{the} \textbf{probability} of different beliefs given your sensory input.
So let me give you an intuitive example.
Imagine you ' re learning to play tennis and you want to decide where \textbf{the} ball is going to bounce as it comes over \textbf{the} net towards you.
There are two sources of information Bayes ' rule tells you.
There ' s sensory evidence — you can use visual information auditory information, and that might tell you it ' s going to land in that red spot.
But you know that your senses are not perfect, and therefore there ' s some variability of where it ' s going to land shown by that cloud of red, representing numbers between 0. 5 and maybe 0. 1.
That information is available in \textbf{the} current shot, but there ' s another source of information not available on \textbf{the} current shot, but only available by repeated experience in \textbf{the} game of tennis, and that ' s that \textbf{the} ball doesn ' t bounce with equal \textbf{probability} over \textbf{the} court during \textbf{the} match.
If you ' re playing against a very good opponent, they may distribute it in that green area, which is \textbf{the} prior distribution, making it hard for you to return.
Now both these sources of information carry important information.
And what Bayes ' rule says is that I should multiply \textbf{the} numbers on \textbf{the} red by \textbf{the} numbers on \textbf{the} green to get \textbf{the} numbers of \textbf{the} \textbf{yellow}, which have \textbf{the} ellipses, and that ' s my belief.
So it ' s \textbf{the} optimal way of combining information.
Now I wouldn ' t tell you all this if it wasn ' t that a few years ago, we showed this is exactly what people do when they learn new movement skills.
And what it means is we really are Bayesian inference machines.
As we go around, we learn about statistics of \textbf{the} world and lay that down, but we also learn about how noisy our own sensory apparatus is, and then combine those in a real Bayesian way.
Now a key part to \textbf{the} Bayesian is this part of \textbf{the} formula.
And what this part really says is I have to predict \textbf{the} \textbf{probability} of different sensory feedbacks given my beliefs.
So that really means I have to make predictions of \textbf{the} future.
And I want to convince you \textbf{the} brain does make predictions of \textbf{the} sensory feedback it ' s going to get.
And moreover, it profoundly changes your perceptions by what you do.
And to do that, I ' ll tell you about how \textbf{the} brain deals with sensory input.
So you send a command out, you get sensory feedback back, and that transformation is governed by \textbf{the} physics of your body and your sensory apparatus.
But you can imagine looking inside \textbf{the} brain.
And here ' s inside \textbf{the} brain.
You might have a little predictor, a neural simulator, of \textbf{the} physics of your body and your senses.
So as you send a movement command down, you tap a \textbf{copy} of that off and run it into your neural simulator to anticipate \textbf{the} sensory consequences of your actions.
So as I shake this ketchup bottle, I get some true sensory feedback as \textbf{the} function of time in \textbf{the} bottom row.
And if I ' ve got a good predictor, it predicts \textbf{the} same thing.
Well why would I bother doing that?
I ' m going to get \textbf{the} same feedback anyway.
Well there ' s good reasons.
Imagine, as I shake \textbf{the} ketchup bottle, someone very kindly comes up to me and taps it on \textbf{the} back for me.
Now I get an extra source of sensory information due to that external act.
So I get two sources.
I get you tapping on it, and I get me shaking it, but from my senses ' point of view, that is combined together into one source of information.
Now there ' s good reason to believe that you would want to be able to distinguish external events from internal events.
Because external events are actually much more behaviorally relevant than feeling everything that ' s going on inside my body.
So one way to reconstruct that is to compare \textbf{the} prediction — which is only based on your movement commands — with \textbf{the} reality.
Any discrepancy should hopefully be external.
So as I go around \textbf{the} world, I ' m making predictions of what I should get, subtracting them off.
Everything left over is external to me.
What evidence is there for this?
Well there ' s one very clear example where a sensation generated by myself feels very different then if generated by another person.
And so we decided \textbf{the} most obvious place to start was with tickling.
It ' s been known for a long time, you can ' t tickle yourself as well as other people can.
But it hasn ' t really been shown, it ' s because you have a neural simulator, simulating your own body and subtracting off that sense.
So we can bring \textbf{the} experiments of \textbf{the} 21st century by applying robotic technologies to this problem.
And in effect, what we have is some sort of stick in one hand attached to a robot, and they ' re going to move that back and forward.
And then we ' re going to track that with a \textbf{computer} and use it to control another robot, which is going to tickle their palm with another stick.
And then we ' re going to ask them to rate a bunch of things including ticklishness.
I ' ll show you just one part of our study.
And here I ' ve taken away \textbf{the} robots, but basically people move with their right arm sinusoidally back and forward.
And we replay that to \textbf{the} other hand with a time delay.
Either no time delay, in which case light would just tickle your palm, or with a time delay of two- tenths of three-tenths of a second.
So \textbf{the} important point here is \textbf{the} right hand always does \textbf{the} same things — sinusoidal movement. \textbf{The} left hand always is \textbf{the} same and puts sinusoidal tickle.
All we ' re playing with is a tempo causality.
And as we go from naught to 0. 1 second, it becomes more ticklish.
As you go from 0. 1 to 0. 2, it becomes more ticklish at \textbf{the} end.
And by 0. 2 of a second, it ' s equivalently ticklish to \textbf{the} robot that just tickled you without you doing anything.
So whatever is responsible for this cancellation is extremely tightly coupled with tempo causality.
And based on this illustration, we really convinced ourselves in \textbf{the} field that \textbf{the} brain ' s making precise predictions and subtracting them off from \textbf{the} sensations.
Now I have to admit, these are \textbf{the} worst studies my lab has ever run.
Because \textbf{the} tickle sensation on \textbf{the} palm comes and goes, you need large numbers of subjects with these stars making them significant.
So we were looking for a much more objective way to assess this phenomena.
And in \textbf{the} intervening years I had two daughters.
And one thing you notice about children in backseats of cars on long journeys, they get into fights — which started with one of them doing something to \textbf{the} other, \textbf{the} other retaliating.
It quickly escalates.
And children tend to get into fights which escalate in terms of force.
Now when I screamed at my children to stop, sometimes they would both say to me \textbf{the} other person hit them harder.
Now I happen to know my children don ' t lie, so I thought, as a neuroscientist, it was important how I could explain how they were telling inconsistent truths.
And we hypothesize based on \textbf{the} tickling study that when one child hits another, they generate \textbf{the} movement command.
They predict \textbf{the} sensory consequences and subtract it off.
So they actually think they ' ve hit \textbf{the} person less hard than they have — rather like \textbf{the} tickling.
Whereas \textbf{the} passive recipient doesn ' t make \textbf{the} prediction, feels \textbf{the} full blow.
So if they retaliate with \textbf{the} same force, \textbf{the} first person will think it ' s been escalated.
So we decided to test this in \textbf{the} lab.
Now we don ' t work with children, we don ' t work with hitting, but \textbf{the} concept is \textbf{identical}.
We bring in two adults.
We tell them they ' re going to play a game.
And so here ' s player one and player two sitting opposite to each other.
And \textbf{the} game is very simple.
We started with a motor with a little lever, a little force transfuser.
And we use this motor to apply force down to player one ' s fingers for three seconds and then it stops.
And that player ' s been told, remember \textbf{the} experience of that force and use your other finger to apply \textbf{the} same force down to \textbf{the} other subject ' s finger through a force transfuser — and they do that.
And player two ' s been told, remember \textbf{the} experience of that force.
Use your other hand to apply \textbf{the} force back down.
And so they take it in turns to apply \textbf{the} force they ' ve just experienced back and forward.
But critically, they ' re briefed about \textbf{the} rules of \textbf{the} game in separate rooms.
So they don ' t know \textbf{the} rules \textbf{the} other person ' s playing by.
And what we ' ve measured is \textbf{the} force as a function of terms.
And if we look at what we start with, a quarter of a Newton there, a number of turns, perfect would be that red line.
And what we see in all pairs of subjects is this — a 70 percent escalation in force on each go.
So it really suggests, when you ' re doing this — based on this study and others we ' ve done — that \textbf{the} brain is canceling \textbf{the} sensory consequences and underestimating \textbf{the} force it ' s producing.
So it re-shows \textbf{the} brain makes predictions and fundamentally changes \textbf{the} precepts.
So we ' ve made inferences, we ' ve done predictions, now we have to generate actions.
And what Bayes ' rule says is, given my beliefs, \textbf{the} action should in some sense be optimal.
But we ' ve got a problem.
Tasks are symbolic — I want to \textbf{drink}, I want to dance — but \textbf{the} movement system has to contract 600 muscles in a particular sequence.
And there ' s a big gap between \textbf{the} task and \textbf{the} movement system.
So it could be bridged in infinitely many different ways.
So think about just a point to point movement.
I could choose these two paths out of an infinite number of paths.
Having chosen a particular path, I can hold my hand on that path as infinitely many different joint \textbf{configurations}.
And I can hold my arm in a particular joint \textbf{configuration} either very stiff or very relaxed.
So I have a huge amount of choice to make.
Now it turns out, we are extremely stereotypical.
We all move \textbf{the} same way pretty much.
And so it turns out we ' re so stereotypical, our brains have got dedicated neural circuitry to decode this stereotyping.
So if I take some dots and set them in motion with biological motion, your brain ' s circuitry would understand instantly what ' s going on.
Now this is a bunch of dots moving.
You will know what this person is doing, whether happy, sad, old, young — a huge amount of information.
If these dots were cars going on a racing circuit, you would have absolutely no idea what ' s going on.
So why is it that we move \textbf{the} particular ways we do?
Well let ' s think about what really happens.
Maybe we don ' t all quite move \textbf{the} same way.
Maybe there ' s variation in \textbf{the} population.
And maybe those who move better than others have got more chance of getting their children into \textbf{the} next generation.
So in \textbf{evolutionary} scales, movements get better.
And perhaps in life, movements get better through learning.
So what is it about a movement which is good or bad?
Imagine I want to intercept this ball.
Here are two possible paths to that ball.
Well if I choose \textbf{the} left-hand path, I can work out \textbf{the} forces required in one of my muscles as a function of time.
But there ' s noise added to this.
So what I actually get, based on this lovely, smooth, desired force, is a very noisy version.
So if I pick \textbf{the} same command through many times, I will get a different noisy version each time, because noise changes each time.
So what I can show you here is how \textbf{the} variability of \textbf{the} movement will \textbf{evolve} if I choose that way.
If I choose a different way of moving — on \textbf{the} right for example — then I ' ll have a different command, different noise, playing through a noisy system, very complicated.
All we can be sure of is \textbf{the} variability will be different.
If I move in this particular way, I end up with a smaller variability across many movements.
So if I have to choose between those two, I would choose \textbf{the} right one because it ' s less variable.
And \textbf{the} fundamental idea is you want to plan your movements so as to minimize \textbf{the} negative consequence of \textbf{the} noise.
And one intuition to get is actually \textbf{the} amount of noise or variability I show here gets bigger as \textbf{the} force gets bigger.
So you want to avoid big forces as one principle.
So we ' ve shown that using this, we can explain a huge amount of data — that exactly people are going about their lives planning movements so as to minimize negative consequences of noise.
So I hope I ' ve convinced you \textbf{the} brain is there and \textbf{evolved} to control movement.
And it ' s an intellectual challenge to understand how we do that.
But it ' s also relevant for disease and rehabilitation.
There are many diseases which effect movement.
And hopefully if we understand how we control movement, we can apply that to robotic technology.
And finally, I want to remind you, when you see animals do what look like very simple tasks, \textbf{the} actual complexity of what is going on inside their brain is really quite dramatic.
Thank you very much.
Chris Anderson : Quick question for you, Dan.
So you ' re a movement — — chauvinist.
Does that mean that you think that \textbf{the} other things we think our brains are about — \textbf{the} dreaming, \textbf{the} yearning, \textbf{the} falling in love and all these things — are a kind of side show, an accident?
DW : No, no, actually I think they ' re all important to drive \textbf{the} right movement behavior to get \textbf{reproduction} in \textbf{the} end.
So I think people who study sensation or memory without realizing why you ' re laying down memories of childhood. \textbf{The} fact that we forget most of our childhood, for example, is probably fine, because it doesn ' t effect our movements later in life.
You only need to store things which are really going to effect movement.
CA : So you think that people thinking about \textbf{the} brain, and consciousness generally, could get real insight by saying, where does movement play in this game?
DW : So people have found out for example that studying vision in \textbf{the} absence of realizing why you have vision is a mistake.
You have to study vision with \textbf{the} realization of how \textbf{the} movement system is going to use vision.
And it uses it very differently once you think about it that way.
CA : Well that was quite fascinating.
Thank you very much indeed.

% matched lemmas: computer, configuration, copy, drink, evolutionary, evolve, identical, learning, manipulate, planet, probability, reproduction, rock, sea, soup, stack, the, year-old, yellow
\end{textsample}
